% sec-08: Document 08, the building code standards for the La Jolla facility.
\docpart{Building Code}{Building Code Standards for a PDAC LLM Oncology Clinical
Trial Facility}

\section{Scope, Occupancy, and Applicable Codes}

\draftnote{Fix the occupancy classification and name every code that governs:
the California Building Standards Code, Title 24, NFPA 99 for health care
facilities, and NFPA 101 for life safety. State which parts of the facility are
business occupancy and which are ambulatory care, because the two carry
different sprinkler and egress requirements and the difference decides the
tenant improvement cost in document 14 \sect{}6.}

\section{Site and Structural Provisions}

\begin{mvtable}
\begin{tabularx}{\textwidth}{>{\raggedright\arraybackslash}p{4.0cm} >{\raggedright\arraybackslash}p{3.0cm} >{\raggedright\arraybackslash}X}
\headrow \thead{Element} & \thead{Requirement} & \thead{Why this trial needs it}\\
\midrule
Procedure suite floor loading & to be fixed & to be fixed \\
Vibration limit at the robot base & to be fixed & to be fixed \\
\end{tabularx}
\end{mvtable}
\tabcapl{tab:structural}{Structural provisions. Stage 2 fixes each value and
states the consequence of not meeting it, which for a robot base is a
registration error rather than a structural failure.}

\draftnote{Derive the robotic arm base and vibration provisions from the device
envelope described in papers 9 and 10 of the deck table at
\mvfile{funding/move-in/inputs/READMES/README-LLM-Pancreatic-Oncology-Clinical-Trial-System-Large-Documents-Funding-and-AI-Peer-Review.md}.
Cite IEC 80601-2-77 for robotically assisted surgical equipment.}

\section{Mechanical, Electrical, and Plumbing}

\draftnote{Air changes per hour and pressure relationships for the procedure
suites, the pharmacy compounding area, and the specimen laboratory; essential
electrical system branches; and the plumbing provisions for the compounding
area. Take the pharmacy provisions from USP General Chapters 797 and 800, and
keep them identical to document 13 \sect{}5.}

\section{Medical Gas, Shielding, and Specialty Systems}

\draftnote{Medical gas outlets per suite, imaging shielding where a fluoroscopic
or computed tomography unit is present, and the specialty systems the protocol
requires. If the protocol does not require an imaging modality on site, say so
and state where the imaging is performed instead, citing
\mvfile{trial-protocol/}.}

\section{Technology Infrastructure and the Machine Room}

\draftnote{The on-premises model node is the reason this document differs from
an ordinary ambulatory surgical center building code. Fix the machine room
provisions: power, cooling, physical access, network segmentation, and the
uninterruptible supply runtime. Keep the segmentation identical to document 09
\sect{}2 and document 11 \sect{}5. Take the on-premises requirement from
\mvfile{trial-protocol/} rather than asserting it.}

\section{Commissioning and Certificate of Occupancy}

\draftnote{The commissioning sequence as a gated table, ending at the
certificate of occupancy, and the handoff to the activation checklist in
document 12 \sect{}2. The two must share a boundary: the last commissioning gate
is the first activation gate.}
